%%% =======================================================================
%%% ISLS Extension Phase 5.1 v1.0.0
%%% C24: Recursive Decomposition and Hierarchical Composition Engine
%%% Author: Sebastian Klemm
%%% Date: 16 March 2026
%%% Status: Implementation Specification for Claude Code
%%% =======================================================================
\documentclass[11pt,a4paper,twoside]{article}

\usepackage[utf8]{inputenc}
\usepackage{amsmath,amssymb,amsthm,mathtools}
\usepackage{algorithm,algorithmic}
\usepackage{booktabs,longtable,tabularx}
\usepackage{enumitem}
\usepackage{geometry}
\usepackage{hyperref}
\usepackage{cleveref}
\usepackage{xcolor}
\usepackage{fancyhdr}
\usepackage{listings}
\usepackage{tikz}
\usetikzlibrary{arrows.meta,positioning,shapes.geometric,fit,calc,trees}

\geometry{margin=2.5cm,top=3cm,bottom=3cm,headheight=14pt}

\theoremstyle{definition}
\newtheorem{definition}{Definition}[section]
\newtheorem{axiom}{Axiom}[section]
\newtheorem{requirement}{Requirement}[section]
\newtheorem{invariant}{Invariant}[section]

\theoremstyle{plain}
\newtheorem{theorem}{Theorem}[section]
\newtheorem{proposition}{Proposition}[section]
\newtheorem{lemma}{Lemma}[section]
\newtheorem{corollary}{Corollary}[section]

\theoremstyle{remark}
\newtheorem{remark}{Remark}[section]

\newcommand{\ISLS}{\textsf{ISLS}}
\newcommand{\FORGE}{\textsf{FORGE}}
\newcommand{\PMHD}{\textsf{PMHD}}
\newcommand{\RR}{\mathbb{R}}
\newcommand{\NN}{\mathbb{N}}
\newcommand{\calH}{\mathcal{H}}
\newcommand{\calC}{\mathcal{C}}
\newcommand{\calD}{\mathcal{D}}
\newcommand{\calS}{\mathcal{S}}
\newcommand{\calT}{\mathcal{T}}
\newcommand{\calI}{\mathcal{I}}
\newcommand{\calU}{\mathcal{U}}
\DeclareMathOperator{\Dig}{Dig}
\DeclareMathOperator{\Canon}{Canon}
\DeclareMathOperator{\Decomp}{Decomp}
\DeclareMathOperator{\Compose}{Compose}
\DeclareMathOperator{\Resolve}{Resolve}
\newcommand{\norm}[1]{\left\lVert #1 \right\rVert}

\lstset{
  basicstyle=\ttfamily\small,
  breaklines=true,
  frame=single,
  backgroundcolor=\color{gray!5},
  rulecolor=\color{gray!30},
}

\pagestyle{fancy}
\fancyhf{}
\fancyhead[LE,RO]{\ISLS{} Extension Phase 5.1 v1.0.0}
\fancyhead[RE,LO]{Sebastian Klemm}
\fancyfoot[C]{\thepage}

\hypersetup{
  colorlinks=true,
  linkcolor=blue!60!black,
  pdftitle={ISLS Extension Phase 5.1 — Recursive Composition},
  pdfauthor={Sebastian Klemm},
}

\begin{document}

\begin{titlepage}
\centering
\vspace*{2cm}
{\Huge\bfseries \ISLS{} Extension Phase~5.1\\[0.3cm]
v1.0.0}\\[1.5cm]
{\Large C24: Recursive Decomposition and\\
Hierarchical Composition Engine}\\[1cm]
{\large Hypercubes of Hypercubes:\\
From Intent to System through\\
Atomic Crystallization and Upward Composition}\\[2cm]
{\large Sebastian Klemm}\\[0.5cm]
{\large 16 March 2026}\\[1.5cm]

\begin{abstract}
\noindent
This document specifies \texttt{isls-compose} (C24), a recursive
decomposition and hierarchical composition engine that extends the
\ISLS{} Forge (C21--C23) from producing single artifacts to producing
\emph{systems of artifacts} --- complete software components, services,
or multi-module architectures.

The core mechanism: a top-level DecisionSpec is recursively decomposed
into atomic sub-specs. Each atom is independently drilled through PMHD
(C21), crystallized through the 8-gate cascade, and emitted as an
ArtifactIR (C22). The atomic crystals are then composed upward through
interface resolution, dependency validation, and hierarchical
crystallization --- producing a \emph{System Crystal} that contains
the complete artifact tree with full provenance at every level.

Each crystal at every level is a Hypercube Universe (C18). Interfaces
between components are Bridges. The atom-molecule-system hierarchy uses
Scale Ladders. The entire recursive process is deterministic,
replay-verified, and evidence-bound.

The result is a generative engine that can produce arbitrarily complex
structured outputs --- not by rendering templates, but by recursive
adversarial crystallization from intent down to atoms and composition
back up to systems.

\medskip\noindent\textbf{Status:} Implementation Specification.\\
\textbf{Parent:} ISLS v1.0.0 + Extensions Phase 1--5.\\
\textbf{Depends on:} C18 (Scale), C21 (PMHD), C22 (ArtifactIR),
  C23 (Forge).
\end{abstract}
\end{titlepage}

\tableofcontents
\newpage


%%% =====================================================================
\part{Foundations}\label{part:foundations}

\section{The Problem with Flat Synthesis}\label{sec:problem}

Phase~5 produces single artifacts: one DecisionSpec $\to$ one Monolith
$\to$ one ArtifactIR $\to$ one output. This is sufficient for
generating a function, a schema, or a configuration file. It is not
sufficient for generating a \emph{system}: a multi-component
architecture with internal interfaces, dependency ordering, layered
abstractions, and compositional correctness.

A system is not a large artifact. It is a \emph{tree of artifacts}
where each node is independently valid and the tree as a whole
satisfies compositional constraints that no single node can express.

\section{Axioms}\label{sec:axioms}

\begin{axiom}[Recursive Decomposability]\label{ax:decomp}
Any DecisionSpec that describes a system can be decomposed into a
finite tree of sub-specs, where each leaf spec describes an atomic
component that is synthesizable by a single Forge run.
\end{axiom}

\begin{axiom}[Independent Atomic Validity]\label{ax:atomvalid}
Each atomic component must independently pass the 8-gate validation.
Compositional validity does not compensate for atomic invalidity.
\end{axiom}

\begin{axiom}[Compositional Closure]\label{ax:compclosure}
The composition of individually valid components into a system must
itself be validated through the 8-gate cascade. The system crystal is
not merely the union of its parts --- it is a new crystal that proves
the composition is sound.
\end{axiom}

\begin{axiom}[Interface-First Decomposition]\label{ax:interface}
Decomposition is driven by interfaces, not by content. The
decomposition algorithm identifies \emph{what needs to connect to
what} before deciding \emph{what each component does}. Interfaces are
the invariant; implementations are the variables.
\end{axiom}

\begin{axiom}[Scale Correspondence]\label{ax:scalecorr}
The decomposition hierarchy maps directly to the Scale architecture
(C18): atoms are micro-scale Hypercube Universes, modules are
meso-scale clusters, and the system is the macro-scale domain.
Bridges between atoms are component interfaces. Ladders between
levels are the composition operators.
\end{axiom}


\section{Three Levels}\label{sec:levels}

\begin{definition}[Atom]\label{def:atom}
An atom is a leaf-level component: a single function, type, endpoint,
schema definition, or configuration block. An atom has no internal
sub-components. It is the smallest unit that can be independently
drilled, crystallized, and validated.
\end{definition}

\begin{definition}[Molecule]\label{def:molecule}
A molecule is a composition of atoms that together form a coherent
module: a service, a library, a layer. A molecule has internal
interfaces (between its atoms) and external interfaces (exposed to
other molecules or to the system).
\end{definition}

\begin{definition}[System]\label{def:system}
A system is a composition of molecules that together form a complete
deliverable: an application, a service mesh, a platform component.
The system crystal is the root of the composition tree.
\end{definition}

\begin{remark}
The three levels are not fixed. For very complex systems, the hierarchy
may have more than three levels (atoms $\to$ molecules $\to$ assemblies
$\to$ systems $\to$ platforms). The decomposition engine supports
arbitrary depth. Three levels is the minimum useful hierarchy.
\end{remark}


%%% =====================================================================
\part{Decomposition}\label{part:decomp}

\section{The Decomposition Tree}\label{sec:decomptree}

\begin{definition}[CompositionTree]\label{def:comptree}
A CompositionTree $\calT$ is a rooted tree where:
\begin{itemize}[nosep]
  \item The root is the top-level DecisionSpec $\calD_0$.
  \item Each internal node is a DecisionSpec $\calD_i$ with children
        $\{\calD_{i,1}, \ldots, \calD_{i,k_i}\}$.
  \item Each leaf is an atomic DecisionSpec that is directly
        forge-able by C23.
  \item Each edge carries an InterfaceContract specifying what the
        child provides to or requires from its parent and siblings.
\end{itemize}
\end{definition}

\begin{definition}[InterfaceContract]\label{def:contract}
An InterfaceContract $\calI$ between two nodes specifies:
\[
  \calI = \langle \mathtt{provider}, \mathtt{consumer},
          \mathtt{provides}, \mathtt{requires},
          \mathtt{protocol}, \mathtt{direction} \rangle
\]
where:
\begin{itemize}[nosep]
  \item \texttt{provides}: list of named capabilities with type
        signatures (e.g., \texttt{"authenticate(token) -> User"}),
  \item \texttt{requires}: list of named dependencies with type
        signatures,
  \item \texttt{protocol}: interaction pattern
        (\texttt{sync\_call}, \texttt{async\_message},
        \texttt{shared\_type}, \texttt{event\_stream}),
  \item \texttt{direction}: $\texttt{Unidirectional}$ or
        $\texttt{Bidirectional}$.
\end{itemize}
\end{definition}


\section{Decomposition Algorithm}\label{sec:decompalgo}

\begin{algorithm}[H]
\caption{\textsc{Decompose}: Recursive spec decomposition}
\label{alg:decompose}
\begin{algorithmic}[1]
\REQUIRE Top-level DecisionSpec $\calD_0$, decomposition config $\theta$
\ENSURE CompositionTree $\calT$
\STATE $\calT \gets \textsc{NewTree}(\calD_0)$
\STATE $\textsc{DecomposeNode}(\calT, \calD_0, 0)$
\RETURN $\calT$
\end{algorithmic}
\end{algorithm}

\begin{algorithm}[H]
\caption{\textsc{DecomposeNode}: Recursive node decomposition}
\label{alg:decompnode}
\begin{algorithmic}[1]
\REQUIRE Tree $\calT$, spec $\calD$, depth $d$
\IF{$\textsc{IsAtomic}(\calD, \theta) \;\vee\; d \ge \theta.\mathrm{max\_depth}$}
  \STATE Mark $\calD$ as leaf in $\calT$
  \RETURN
\ENDIF
\STATE $\{(\calD_1, \calI_1), \ldots, (\calD_k, \calI_k)\} \gets
  \textsc{IdentifyComponents}(\calD, \theta)$
\FOR{$j = 1$ \TO $k$}
  \STATE Add $\calD_j$ as child of $\calD$ in $\calT$ with edge $\calI_j$
  \STATE $\textsc{DecomposeNode}(\calT, \calD_j, d + 1)$
\ENDFOR
\end{algorithmic}
\end{algorithm}

\begin{definition}[Atomicity Test]\label{def:atomtest}
A DecisionSpec is atomic if and only if:
\begin{enumerate}[nosep,label=(\roman*)]
  \item it describes a single responsibility (one primary capability),
  \item its estimated component count is $\le \theta_{\mathrm{atom\_max}}$
        (default: 5 components in the ArtifactIR),
  \item it has no natural internal interface boundaries,
  \item or the depth limit has been reached.
\end{enumerate}
\end{definition}

\begin{definition}[Component Identification]\label{def:compident}
Component identification analyzes a non-atomic DecisionSpec and splits
it along interface boundaries:
\begin{enumerate}[nosep,label=(\roman*)]
  \item \textbf{Capability clustering}: group related capabilities
        (e.g., all auth-related capabilities form one component).
  \item \textbf{Dependency analysis}: identify which capabilities
        depend on which others; cut at minimal-dependency boundaries.
  \item \textbf{Interface extraction}: for each cut, define the
        InterfaceContract (what crosses the boundary).
  \item \textbf{Sub-spec generation}: create a new DecisionSpec for
        each cluster, inheriting relevant goals and constraints from
        the parent, plus new constraints from the interface contracts.
\end{enumerate}
The identification is deterministic: same spec and config produce
identical decomposition.
\end{definition}


\section{Decomposition Strategies}\label{sec:decompstrat}

\begin{definition}[Decomposition Strategies]\label{def:strategies}
The decomposition engine supports multiple strategies for splitting
a spec into sub-specs:
\begin{enumerate}[nosep,label=(\roman*)]
  \item \textbf{LayerDecomp}: split by architectural layers
        (e.g., data layer, business logic, API layer, presentation).
  \item \textbf{FeatureDecomp}: split by features / capabilities
        (each feature becomes a component).
  \item \textbf{DomainDecomp}: split by domain boundaries
        (Domain-Driven Design-style bounded contexts).
  \item \textbf{PipelineDecomp}: split by processing stages
        (input $\to$ transform $\to$ validate $\to$ output).
  \item \textbf{HybridDecomp}: combine strategies per level
        (e.g., LayerDecomp at system level, FeatureDecomp at module
        level).
  \item \textbf{CustomDecomp}: user-provided decomposition function.
\end{enumerate}
All strategies are deterministic.
\end{definition}


%%% =====================================================================
\part{Per-Atom Forging}\label{part:atomforge}

\section{Atomic Forge Pipeline}\label{sec:atomforge}

Each leaf in the CompositionTree runs through the standard Phase~5
forge pipeline independently:

\begin{enumerate}[nosep]
  \item PMHD Drill on the atomic DecisionSpec $\to$ Monolith
  \item ArtifactIR from Monolith
  \item Matrix interpretation (domain-specific)
  \item Synthesis $\to$ raw output
  \item Evaluation against constraints
  \item 8-gate validation $\to$ Atomic Crystal $C_{\mathrm{atom}}$
\end{enumerate}

\begin{invariant}[Atomic Independence]\label{inv:atomind}
Each atomic forge run is independent: it does not read the output of
sibling atoms. It knows about siblings only through the
InterfaceContracts attached to its DecisionSpec. This ensures that
each atom can be forged in parallel and that atomic validity is
self-contained.
\end{invariant}

\begin{requirement}[Interface Constraint Injection]\label{req:iconstinject}
When an atom's DecisionSpec is created by decomposition, the
InterfaceContracts from its parent and siblings are injected as
additional hard constraints:
\begin{itemize}[nosep]
  \item \texttt{provides} entries become ``must export'' constraints.
  \item \texttt{requires} entries become ``must import'' constraints.
  \item Type signatures from contracts become compatibility constraints.
\end{itemize}
This ensures that independently forged atoms produce compatible
interfaces.
\end{requirement}


\section{Parallel Forging}\label{sec:parallel}

\begin{definition}[Forge Schedule]\label{def:forgesched}
The forge schedule determines the order of atomic forge runs. Since
atoms are independent (by \cref{inv:atomind}), all atoms at the same
depth level can be forged in parallel. The schedule is a sequence of
\emph{waves}:
\[
  W_d = \{C_{\mathrm{atom}} : \mathrm{depth}(C) = d_{\max}\}
\]
Wave $d_{\max}$ (deepest leaves) is forged first, then
$d_{\max} - 1$, up to the root. Within each wave, all atoms are
independent and may be forged concurrently.
\end{definition}

\begin{requirement}[Sequential Fallback]\label{req:seqfallback}
On single-core systems, atoms within a wave are forged sequentially
in deterministic order (sorted by sub-spec content address). This
preserves determinism regardless of parallelism.
\end{requirement}


%%% =====================================================================
\part{Composition}\label{part:compose}

\section{Interface Resolution}\label{sec:resolve}

\begin{definition}[Interface Resolution]\label{def:resolve}
After all atoms in a molecule are forged, their interfaces are
resolved:
\[
  \Resolve(\{C_1, \ldots, C_k\}, \{\calI_1, \ldots, \calI_m\})
  = (\Gamma, \Delta)
\]
where $\Gamma$ is the set of \emph{satisfied} interface bindings and
$\Delta$ is the set of \emph{unsatisfied} requirements.
\end{definition}

\begin{definition}[Interface Binding]\label{def:binding}
A binding $(p, c, \calI)$ is satisfied when:
\begin{enumerate}[nosep,label=(\roman*)]
  \item Component $p$ declares a \texttt{provides} entry matching
        $\calI.\mathtt{provides}$.
  \item Component $c$ declares a \texttt{requires} entry matching
        $\calI.\mathtt{requires}$.
  \item The type signatures are compatible (structural subtyping: the
        provider's signature is at least as general as the consumer's
        expectation).
  \item The protocol matches.
\end{enumerate}
\end{definition}

\begin{definition}[Resolution Failure Modes]\label{def:resfail}
\begin{enumerate}[nosep,label=(\roman*)]
  \item \textbf{Unsatisfied requirement}: a consumer needs something
        no provider offers. The composition cannot proceed.
  \item \textbf{Type mismatch}: provider and consumer have incompatible
        type signatures. The offending atom must be re-forged with
        corrected constraints.
  \item \textbf{Circular dependency}: the dependency graph among atoms
        contains a cycle. The decomposition must be revised.
  \item \textbf{Ambiguous binding}: multiple providers match a single
        requirement. Resolution uses deterministic tie-breaking
        (provider with highest quality score, then by content address).
\end{enumerate}
\end{definition}

\begin{requirement}[Resolution is Deterministic]\label{req:resdet}
Interface resolution \textbf{MUST} produce identical bindings under
identical inputs. All tie-breaking rules are deterministic and
specified in the Run Descriptor.
\end{requirement}


\section{Upward Crystallization}\label{sec:upcrystal}

\begin{definition}[Molecule Crystal]\label{def:molcrystal}
A molecule crystal $C_{\mathrm{mol}}$ is produced by composing its
atomic crystals:
\[
  C_{\mathrm{mol}} = \Compose(\{C_{\mathrm{atom},1}, \ldots,
  C_{\mathrm{atom},k}\}, \Gamma, \calD_{\mathrm{mol}})
\]

The molecule crystal contains:
\begin{itemize}[nosep]
  \item \texttt{sub\_crystal\_ids}: content addresses of all atomic
        crystals,
  \item \texttt{interface\_bindings}: the resolved binding set $\Gamma$,
  \item \texttt{composition\_proof}: evidence that all interfaces are
        satisfied and the dependency graph is acyclic,
  \item \texttt{constraint\_program}: the union of atomic constraint
        programs plus compositional constraints,
  \item its own \texttt{topology\_signature}: the dependency graph
        topology of its atoms (using Betti numbers, spectral gap,
        etc.\ from C16).
\end{itemize}
\end{definition}

\begin{definition}[Composition Validation]\label{def:compval}
The molecule crystal passes through a dedicated composition
validation before the standard 8-gate cascade:
\begin{enumerate}[nosep,label=(CV\arabic*)]
  \item \textbf{All atoms valid}: every sub-crystal passes 8/8 checks.
  \item \textbf{No unsatisfied interfaces}: $\Delta = \emptyset$.
  \item \textbf{No type mismatches}: all bindings type-check.
  \item \textbf{Acyclic dependencies}: topological sort succeeds.
  \item \textbf{Coverage}: every capability from the parent
        DecisionSpec is provided by at least one atom.
  \item \textbf{No orphans}: every atom provides at least one
        capability that is consumed.
\end{enumerate}
After CV1--CV6 pass, the molecule crystal enters the standard
8-gate cascade.
\end{definition}

\begin{definition}[System Crystal]\label{def:syscrystal}
The system crystal $C_{\mathrm{sys}}$ is produced by composing
molecule crystals. Its structure is identical to a molecule crystal
but at one level higher:
\begin{itemize}[nosep]
  \item \texttt{sub\_crystal\_ids}: content addresses of molecule
        crystals,
  \item each molecule crystal itself references its atomic crystals,
  \item the system crystal's \texttt{topology\_signature} describes
        the inter-module dependency graph.
\end{itemize}
The system crystal passes CV1--CV6 on its molecule children, then
the 8-gate cascade.
\end{definition}

\begin{theorem}[Hierarchical Validity]\label{thm:hiervalid}
If the system crystal passes CV1--CV6 and the 8-gate cascade, then
every crystal at every level of the hierarchy is independently valid:
every atom passes 8/8, every molecule passes CV1--CV6 + 8/8, and the
system passes CV1--CV6 + 8/8. The validity is not inherited downward
--- it is independently established at each level and verified
compositionally at each parent level.
\end{theorem}

\begin{proof}
By construction: CV1 at the system level requires that all molecule
sub-crystals pass 8/8. CV1 at the molecule level requires that all
atomic sub-crystals pass 8/8. The 8-gate cascade at each level
operates on the crystal's own fields independently. Therefore
validity at level $n$ implies validity at all levels $< n$, and
each level's validity is independently checkable.
\end{proof}


\section{Re-Forging on Failure}\label{sec:reforge}

\begin{definition}[Composition Repair]\label{def:repair}
When interface resolution fails (type mismatch or unsatisfied
requirement), the composition engine can attempt repair:
\begin{enumerate}[nosep,label=(\roman*)]
  \item \textbf{Constraint refinement}: add the failing interface
        requirement as an explicit constraint to the offending atom's
        DecisionSpec and re-forge.
  \item \textbf{Adapter synthesis}: forge a new ``glue'' atom whose
        sole purpose is to bridge the type mismatch between two
        incompatible atoms.
  \item \textbf{Decomposition revision}: re-run decomposition with
        different boundaries if multiple atoms fail to compose.
\end{enumerate}
\end{definition}

\begin{requirement}[Repair Budget]\label{req:repairbudget}
The repair loop \textbf{MUST} be bounded by a configurable budget
(max re-forge attempts per atom, max adapter atoms, max decomposition
revisions). When the budget is exhausted, the composition fails with
a detailed diagnostic listing all unresolved interfaces.
\end{requirement}


%%% =====================================================================
\part{Synthesized Output}\label{part:output}

\section{System Artifact Assembly}\label{sec:assembly}

\begin{definition}[SystemArtifact]\label{def:sysartifact}
The final output of a composition run is a SystemArtifact:
\begin{lstlisting}
pub struct SystemArtifact {
    pub system_crystal: SemanticCrystal,
    pub tree: CompositionTree,
    pub atoms: Vec<AtomArtifact>,
    pub molecules: Vec<MoleculeArtifact>,
    pub bindings: Vec<InterfaceBinding>,
    pub manifest: ExecutionManifest,
    pub total_components: usize,
    pub total_interfaces: usize,
    pub depth: usize,
}

pub struct AtomArtifact {
    pub crystal: SemanticCrystal,
    pub ir: ArtifactIR,
    pub synthesis: SynthesisOutput,
    pub file_path: String,
}

pub struct MoleculeArtifact {
    pub crystal: SemanticCrystal,
    pub atom_ids: Vec<Hash256>,
    pub bindings: Vec<InterfaceBinding>,
    pub composition_proof: CompositionProof,
}
\end{lstlisting}
\end{definition}

\begin{definition}[File Layout]\label{def:filelayout}
The emitter produces a file tree reflecting the composition hierarchy:
\begin{lstlisting}
output/
  system.crystal.json          # System crystal
  manifest.json                # Execution manifest
  composition_proof.json       # CV1-CV6 evidence
  modules/
    auth/
      auth.crystal.json        # Molecule crystal
      atoms/
        auth_handler.rs        # Atom artifact
        auth_handler.crystal.json
        token_validator.rs
        token_validator.crystal.json
    database/
      database.crystal.json
      atoms/
        schema.sql
        schema.crystal.json
        connection.rs
        connection.crystal.json
    api/
      api.crystal.json
      atoms/
        routes.rs
        routes.crystal.json
        middleware.rs
        middleware.crystal.json
  interfaces/
    bindings.json              # All resolved interface bindings
    dependency_graph.json      # Topological order
\end{lstlisting}
Every file has a companion crystal. Every crystal is independently
verifiable. The system crystal references all others.
\end{definition}


%%% =====================================================================
\part{Scale Integration}\label{part:scale}

\section{Mapping to Hypercube Universes}\label{sec:scalemap}

\begin{definition}[Composition-Scale Mapping]\label{def:scalemap}
The composition hierarchy maps to the C18 Scale architecture:
\begin{itemize}[nosep]
  \item Each atom is a micro-scale Hypercube Universe. Its vertices are
        the ArtifactIR components. Its edges are component dependencies.
        Its aggregate state is the atom's 5D quality signature.
  \item Each molecule is a meso-scale Hypercube Universe. Its vertices
        are the atom-universes. Its edges are interface bindings
        (Bridges). Its aggregate state is the molecule's composite
        quality.
  \item The system is the macro-scale Hypercube Universe. Its vertices
        are molecule-universes. Its edges are inter-module interfaces.
\end{itemize}
\end{definition}

\begin{definition}[Composition Bridges]\label{def:compbridges}
An interface binding between two atoms becomes a Bridge (C18) between
their Hypercube Universes:
\begin{itemize}[nosep]
  \item Weight = type compatibility score,
  \item Delay = 0 (synchronous composition),
  \item Phase offset = 0,
  \item Activation predicate = always active (bindings are resolved,
        not conditional).
\end{itemize}
\end{definition}

\begin{definition}[Composition Ladders]\label{def:compladders}
The Lift and Projection operators from C18 apply:
\begin{itemize}[nosep]
  \item $\mathrm{Lift}_{\mathrm{atom} \to \mathrm{mol}}$: aggregate
        atom quality metrics into molecule quality.
  \item $\mathrm{Proj}_{\mathrm{mol} \to \mathrm{atom}}$: propagate
        molecule-level constraints down to atoms (for re-forging).
  \item $\mathrm{Lift}_{\mathrm{mol} \to \mathrm{sys}}$: aggregate
        molecule quality into system quality.
\end{itemize}
\end{definition}


%%% =====================================================================
\part{Configuration and API}\label{part:api}

\section{Configuration}\label{sec:config}

\begin{lstlisting}[title={Composition Configuration}]
compose:
  max_depth: 4                # max decomposition levels
  atom_max_components: 5      # atomicity threshold
  decomposition_strategy: "hybrid"
  parallel_forge: false       # true if multi-core
  repair:
    max_reforge_per_atom: 3
    max_adapter_atoms: 5
    max_decomp_revisions: 2
  quality_thresholds:
    min_atom_quality: 0.5     # minimum per-axis quality for atoms
    min_molecule_cv_score: 0.7
    min_system_coverage: 0.9
\end{lstlisting}


\section{Rust API}\label{sec:rustapi}

\begin{lstlisting}
// isls-compose/src/lib.rs

pub struct CompositionTree {
    pub root: TreeNode,
    pub depth: usize,
    pub atom_count: usize,
    pub molecule_count: usize,
}

pub struct TreeNode {
    pub spec: DecisionSpec,
    pub level: CompLevel,
    pub children: Vec<TreeNode>,
    pub interfaces: Vec<InterfaceContract>,
    pub crystal: Option<SemanticCrystal>,
}

pub enum CompLevel { System, Molecule, Atom }

pub struct InterfaceContract {
    pub provider: String,
    pub consumer: String,
    pub provides: Vec<Capability>,
    pub requires: Vec<Capability>,
    pub protocol: Protocol,
    pub direction: Direction,
}

pub struct Capability {
    pub name: String,
    pub signature: String,     // type signature
    pub description: String,
}

pub enum Protocol { SyncCall, AsyncMessage, SharedType, EventStream }
pub enum Direction { Unidirectional, Bidirectional }

pub struct InterfaceBinding {
    pub provider_crystal_id: Hash256,
    pub consumer_crystal_id: Hash256,
    pub contract: InterfaceContract,
    pub compatibility_score: f64,
}

pub struct CompositionProof {
    pub cv_results: [bool; 6],    // CV1-CV6
    pub all_atoms_valid: bool,
    pub unsatisfied: Vec<String>, // empty if success
    pub dependency_order: Vec<Hash256>,
    pub coverage: f64,
}

pub struct CompositionEngine {
    forge: ForgeEngine,
    config: ComposeConfig,
}

impl CompositionEngine {
    pub fn new(forge: ForgeEngine, config: ComposeConfig) -> Self;

    /// Full pipeline: decompose -> forge atoms -> resolve -> compose -> validate
    pub fn compose(&mut self, spec: DecisionSpec) -> Result<SystemArtifact>;

    /// Decompose only (for inspection before forging)
    pub fn decompose(&self, spec: &DecisionSpec) -> Result<CompositionTree>;

    /// Forge all atoms in a tree (after decomposition)
    pub fn forge_atoms(&mut self, tree: &mut CompositionTree) -> Result<Vec<AtomArtifact>>;

    /// Resolve interfaces between forged atoms
    pub fn resolve_interfaces(
        &self, atoms: &[AtomArtifact], contracts: &[InterfaceContract]
    ) -> Result<(Vec<InterfaceBinding>, Vec<String>)>;

    /// Compose atoms into molecules, molecules into system
    pub fn compose_upward(
        &mut self, tree: &mut CompositionTree, atoms: &[AtomArtifact]
    ) -> Result<SystemArtifact>;

    /// Attempt repair on failed composition
    pub fn repair(
        &mut self, tree: &mut CompositionTree, failures: &[String]
    ) -> Result<SystemArtifact>;
}
\end{lstlisting}


\section{CLI Integration}\label{sec:cli}

\begin{lstlisting}
# Compose a full system from a DecisionSpec
isls compose --spec system-intent.json --domain rust --output system/

# Decompose only (inspect the tree before forging)
isls compose --spec system-intent.json --decompose-only

# Compose with specific strategy
isls compose --spec intent.json --strategy layer --max-depth 3

# Compose from existing crystals (re-compose with new interfaces)
isls compose --crystals crystal1.json,crystal2.json --output composed/
\end{lstlisting}


\section{Gateway Integration}\label{sec:gateway}

\begin{lstlisting}[title={New Gateway Endpoints}]
POST /compose          Body: {spec, domain, config}
                       -> SystemArtifact with full crystal tree

POST /compose/decompose  Body: {spec, strategy}
                         -> CompositionTree (preview, no forging)

GET  /compose/status     -> current composition progress
                            (atoms forged / total, molecules composed)
\end{lstlisting}


%%% =====================================================================
\part{Acceptance Tests}\label{part:tests}

\begin{longtable}{l l p{6.5cm}}
\toprule
\textbf{ID} & \textbf{Name} & \textbf{Procedure} \\
\midrule
\endfirsthead \midrule \endhead \bottomrule \endfoot
AT-CO1 & Decomposition determinism & Decompose same spec twice; verify
  identical trees. \\
AT-CO2 & Atom forging & Decompose into 3 atoms; forge each; verify
  3 crystals all pass 8/8. \\
AT-CO3 & Interface resolution & Forge 3 atoms with compatible
  interfaces; resolve; verify all bindings satisfied. \\
AT-CO4 & Type mismatch detection & Forge 2 atoms with incompatible
  types; verify resolution fails with diagnostic. \\
AT-CO5 & Upward crystallization & Compose 3 atoms into molecule;
  verify molecule crystal passes CV1--CV6 + 8/8. \\
AT-CO6 & System crystal & Compose 2 molecules into system; verify
  system crystal references both molecules. \\
AT-CO7 & Hierarchical validity & Validate system crystal; verify
  every sub-crystal at every level independently valid. \\
AT-CO8 & Repair on failure & Create intentional type mismatch;
  verify repair re-forges the atom with corrected constraint. \\
AT-CO9 & Depth limit & Set max\_depth=1; verify decomposition
  produces only atoms (no molecules). \\
AT-CO10 & File layout & Compose a 2-molecule system; verify file
  tree matches expected layout. \\
AT-CO11 & Composition determinism & Compose same spec twice; verify
  identical system crystal ID. \\
AT-CO12 & Scale mapping & Compose system; verify C18 Hypercube
  Universes created for atoms, molecules, and system. \\
\end{longtable}


%%% =====================================================================
\part{Dependency Graph and Handoff}\label{part:handoff}

\section{Dependency Graph}

\begin{lstlisting}
isls-compose (C24)
    depends on:
      isls-pmhd (C21)          -- drill engine
      isls-artifact-ir (C22)   -- IR construction
      isls-forge (C23)         -- synthesis + evaluation
      isls-scale (C18)         -- Hypercube/Bridge/Ladder mapping
      isls-consensus (C5)      -- 8-gate validation
      isls-archive (C7)        -- crystal storage
      isls-manifest (C13)      -- execution manifest
      isls-store (C17)         -- persistence
      isls-types (C1)          -- foundation types
\end{lstlisting}

\section{Test Count}

\begin{tabular}{l r r}
\toprule
\textbf{Phase} & \textbf{New Tests} & \textbf{Cumulative} \\
\midrule
Phase 1--5 & 231 & 231 \\
Phase 5.1 (C24) & 12 & $\ge$ 243 \\
\bottomrule
\end{tabular}

\section{Completion Criteria}

Phase 5.1 is complete when:
\begin{enumerate}[nosep]
  \item Decomposition produces deterministic trees from specs.
  \item Each atom is independently forged and validated (8/8).
  \item Interface resolution detects matches and mismatches.
  \item Molecule crystals pass CV1--CV6 + 8/8.
  \item System crystals reference all sub-crystals.
  \item Repair loop handles type mismatches within budget.
  \item File layout reflects composition hierarchy.
  \item Scale mapping creates Hypercube Universes at all levels.
  \item Total tests $\ge 243$, zero warnings.
  \item \textbf{24 Crates. One binary. Systems from intent.}
\end{enumerate}


\vfill
\begin{center}
\rule{0.5\textwidth}{0.4pt}\\[0.5cm]
{\small Generated for \ISLS{} Extension Phase 5.1 v1.0.0.\\
Atoms crystallize. Molecules compose. Systems emerge.\\
Every level proven. Every interface bound. Every crystal verified.}
\end{center}

\end{document}
